\annexsubsection{\texorpdfstring{Trace of a quadratic form with respect to another \(^{(1)}\)}{Trace of a quadratic form with respect to another (1)}}

In this No., we denote by E a real vector space and by Q, H two positive quadratic forms on E. There exist two symmetric bilinear forms \((x,y)\mapsto(x|y)_{\mathbf{Q}}\) and \((x,y)\mapsto(x|y)_{\mathbf{H}}\) characterized by

\[
\mathrm{Q} (x) = (x | x) _ {\mathrm{Q}}, \quad \mathrm{H} (x) = (x | x) _ {\mathrm{H}}
\]

for all \(x \in E\) .

One calls trace of Q with respect to H and one denotes by \(\operatorname{Tr}\left(\mathbf{Q}/\mathbf{H}\right)\) the positive real number, finite or not, defined as follows:

\begin{enumerate}
\def\labelenumi{\alph{enumi})}
\item
  If there exists an \(x \in \mathbf{E}\) such that \(\mathrm{H}(x) = 0\) and \(\mathrm{Q}(x) \neq 0\) , one sets \(\operatorname{Tr}(\mathrm{Q} / \mathrm{H}) = +\infty\) .
\item
  In the contrary case, \(\operatorname{Tr}(\mathbf{Q}/\mathbf{H})\) is the supremum of the set of numbers of the form \(\sum_{i=1}^{p} \mathbf{Q}(e_i)\) , where \((e_1, \ldots, e_p)\) runs over the set of finite sequences of elements of E orthonormal for H.
\end{enumerate}

Let E be a real Hilbert space and Q a positive quadratic form on E. Set \(\mathrm{H}(x)=\|x\|^{2}\) for all \(x\in E\) ; then H is a positive quadratic form on E. One says that Q is nuclear if \(\mathrm{Tr}\left(\mathrm{Q}/\mathrm{H}\right)\) is finite. For every \(x\in E\) of norm 1, one has \(\mathrm{Q}(x)\leqslant\mathrm{Tr}\left(\mathrm{Q}/\mathrm{H}\right)\) , whence \(Q\leqslant\mathrm{Tr}\left(\mathrm{Q}/\mathrm{H}\right)\cdot\mathrm{H}\) ; in particular, every nuclear form Q is continuous.

Remarks. --- 1) For every subspace F of E, denote by \(Q_{F}\) the restriction of Q to F, and by \(H_{F}\) that of H. Then \(\mathrm{Tr}\left(\mathrm{Q}_{\mathrm{F}}/\mathrm{H}_{\mathrm{F}}\right) \leqslant \mathrm{Tr}\left(\mathrm{Q}/\mathrm{H}\right)\) and \(\mathrm{Tr}\left(\mathrm{Q}/\mathrm{H}\right)\) is the supremum of the numbers \(\mathrm{Tr}\left(\mathrm{Q}_{\mathrm{F}}/\mathrm{H}_{\mathrm{F}}\right)\) for \(F \subset E\) finite-dimensional.

\begin{enumerate}
\def\labelenumi{\arabic{enumi})}
\setcounter{enumi}{1}
\tightlist
\item
  Let \(\mathbf{E}_1\) be a real vector space, \(\mathbf{Q}_1\) and \(\mathbf{H}_1\) two positive quadratic forms on \(\mathbf{E}_1\) , and \(\pi: \mathbf{E} \to \mathbf{E}_1\) a surjective linear mapping. If \(\mathbf{Q} = \mathbf{Q}_1 \circ \pi\) and \(\mathbf{H} = \mathbf{H}_1 \circ \pi\) , then \(\operatorname{Tr}(\mathbf{Q} / \mathbf{H}) = \operatorname{Tr}(\mathbf{Q}_1 / \mathbf{H}_1)\) .
\end{enumerate}

PROPOSITION 1. --- Assume that E is finite-dimensional and H is nondegenerate.

\begin{enumerate}
\def\labelenumi{\alph{enumi})}
\item
  There exists an endomorphism \(u\) of \(\mathbf{E}\) characterized by \((x|y)_{\mathbb{Q}} = (u(x)|y)_{\mathbb{H}}\) for \(x, y\) in \(\mathbf{E}\) .
\item
  \(\operatorname{Tr}(\mathbf{Q} / \mathbf{H}) = \operatorname{Tr}(u)\) .
\item
  \(\operatorname{Tr}(\mathbf{Q} / \mathbf{H}) = \sum_{i=1}^{m} \mathbf{Q}(e_i)\) for every basis \((e_1, \ldots, e_m)\) of E orthonormal for H.
\end{enumerate}

a) follows from the fact that the bilinear form \((x,y)\mapsto (x|y)_{\mathrm{H}}\) is nondegenerate. Every sequence in E orthonormal for H may be completed to a basis of E orthonormal for H. Consequently, \(\operatorname {Tr}(\mathbf{Q} / \mathbf{H})\) is the supremum of the set of numbers of the form \(\sum_{i = 1}^{m}\mathbf{Q}(e_i)\) over all bases \((e_1,\ldots ,e_m)\) of E orthonormal for H. To prove b) and c), it suffices to show that \(\sum_{i = 1}^{m}\mathbf{Q}(e_i) =\) \(\operatorname {Tr}(u)\) for every basis of this kind. Set \(a_{ij} = (u(e_i)|e_j)_{\mathrm{H}} = (e_i|e_j)_{\mathrm{Q}}\) for \(1\leqslant i,j\leqslant m\) ; then \(u(e_i) = \sum_{j = 1}^{m}a_{ij}e_j\) for \(1\leqslant i\leqslant m\) , whence

\[
\operatorname{Tr} (u) = \sum_ {i = 1} ^ {m} a _ {i i} = \sum_ {i = 1} ^ {m} \left(e _ {i} \mid e _ {i}\right) _ {\mathrm{Q}} = \sum_ {i = 1} ^ {m} Q \left(e _ {i}\right).\tag{Q.E.D.}
\]

PROPOSITION 2. --- Assume that E is finite-dimensional. There exist a basis \((e_{1},\ldots,e_{n})\) of E and an integer m with \(0\leqslant m\leqslant n\) such that

\[
\mathrm{H} \left(\sum_ {i = 1} ^ {n} t _ {i} e _ {i}\right) = \sum_ {i = 1} ^ {m} t _ {i} ^ {2}\tag{1}
\]

for \(t_1, \ldots, t_n\) real. If, moreover, the relation \(\mathrm{H}(x) = 0\) implies \(\mathrm{Q}(x) = 0\) for \(x \in \mathbf{E}\) , then \(\operatorname{Tr}(\mathrm{Q}/\mathrm{H}) = \sum_{i=1}^{m} \mathrm{Q}(e_i)\) .

There exists a basis \((e_1', \ldots, e_n')\) of E orthogonal for H. One can assume the basis to be indexed in such a way that \(\mathrm{H}(e_i') > 0\) for \(1 \leqslant i \leqslant m\) and \(\mathrm{H}(e_i') = 0\) for \(m < i \leqslant n\) . Then set \(e_i = e_i' / \mathrm{H}(e_i')^{1/2}\) for \(1 \leqslant i \leqslant m\) and \(e_i = e_i'\) for \(m < i \leqslant n\) ; the relation (1) is satisfied.

Let F be the subspace of E generated by \(e_{m+1}^{\prime},\ldots,e_{n}^{\prime}\) ; it is the set of \(x\in E\) such that \(\mathrm{H}(x)=0\) . Denote by \(\pi\) the canonical mapping of E onto \(E_{1}=E/F\) . Since Q and H are zero on F, there exist two positive quadratic forms \(Q_{1}\) and \(H_{1}\) on \(E_{1}\) such that \(Q=Q_{1}\circ\pi\) and \(H=H_{1}\circ\pi\) . Moreover, \((\pi(e_{1}),\ldots,\pi(e_{m}))\) is a basis of \(E_{1}\) orthonormal for \(H_{1}\) , and therefore \(H_{1}\) is nondegenerate.

By Prop. 1 and Remark 2,

\[
\operatorname{Tr} \left(\mathrm{Q} / \mathrm{H}\right) = \operatorname{Tr} \left(\mathrm{Q} _ {1} / \mathrm{H} _ {1}\right) = \sum_ {i = 1} ^ {m} \mathrm{Q} _ {1} \left(\pi (e _ {i})\right) = \sum_ {i = 1} ^ {m} \mathrm{Q} (e _ {i}).\tag{Q.E.D.}
\]

Remark 3). --- Assume E finite-dimensional and H nondegenerate. Let \((e_1, \ldots, e_n)\) be a basis of E. Set \(q = ((e_i | e_j)_Q)_{1 \leqslant i, j \leqslant n}\) and \(h = ((e_i | e_j)_H)_{1 \leqslant i, j \leqslant n}\) . With notations as in Prop. 1, the matrix of \(u\) for the basis \((e_1, \ldots, e_n)\) of E is equal to \(h^{-1}q\) , whence

\[
\operatorname{Tr} \left(\mathrm{Q} / \mathrm{H}\right) = \operatorname{Tr} \left(h ^ {- 1} q\right) = \operatorname{Tr} \left(q h ^ {- 1}\right).\tag{2}
\]

